% =============================================================
% template.tex  —  Plain TeX
%
% FLAGS (set these near the top of your document):
%
%   \bindingleft   or \bindingright  (which physical side is the spine)
%   \doublesided   or \singlesided   (alternating pages or not)
%   \notesouter    or \notesinner    (which side margin notes appear on)
%
% MARGIN WIDTH (set \marginwidth):
%   absolute:  \marginwidth=3cm
%   relative:  \marginwidth=\dimexpr 0.2\pdfpagewidth \relax
%
% HEAD / FOOT HOOKS (token registers — set in your document):
%   \headinner, \headmiddle, \headouter
%   \footinner, \footmiddle, \footouter
%
%   In double-sided mode the output routine resolves inner/outer
%   to left/right based on \ifodd\pageno and the \binding flag.
%   In single-sided mode inner=left, outer=right throughout.
%
% MARGIN NOTES:
%   \marginnote{REF}{NOTE TEXT}
%   REF is any user-chosen symbol (★, †, 42, …).
%   Deposits REF inline and REF + NOTE in the margin.
%   Placement is approximate; REF is the binding reference.
% =============================================================

% =============================================================
% FLAGS
% Plain TeX has no keyword arguments, so we use \ifXXX booleans.
% =============================================================

\newif\ifbindingleft  \bindinglefttrue   % binding=left  (false = right)
\newif\ifdoublesided  \doublesidedtrue   % sides=double  (false = single)
\newif\ifnotesouter   \notesoutertrue    % notes=outer   (false = inner)

% =============================================================
% PAGE DIMENSIONS
% =============================================================

\pdfpagewidth=210mm
\pdfpageheight=297mm   % A4

\newdimen\marginwidth
\marginwidth=\dimexpr 0.2\pdfpagewidth \relax  % RELATIVE — change to e.g. 3cm for absolute

\newdimen\marginparsep
\marginparsep=0.5em

% Body text width: paperwidth minus both side margins minus margin column
\newdimen\bodywidth
\bodywidth=\dimexpr \pdfpagewidth - 5cm - \marginwidth - \marginparsep \relax
                    % 5cm = inner(2.5cm) + outer(2.5cm) side margins

\hsize=\bodywidth
\vsize=\dimexpr \pdfpageheight - 5cm \relax
                % 5cm = top(2.5cm) + bottom(2.5cm)

% Correct for TeX's built-in 1in offsets:
\hoffset=\dimexpr 2.5cm - 1in \relax
\voffset=\dimexpr 2.5cm - 1in \relax

% =============================================================
% HEAD / FOOT HOOKS
% Six token registers: inner/middle/outer for head and foot.
% "inner" = binding side, "outer" = fore-edge side.
% Default: page number on outer head.
% =============================================================

\newtoks\headinner  \newtoks\headmiddle  \newtoks\headouter
\newtoks\footinner  \newtoks\footmiddle  \newtoks\footouter

\headinner={}
\headmiddle={}
\headouter={\folio}   % \folio = page number in Plain TeX
\footinner={}
\footmiddle={}
\footouter={}

% =============================================================
% OUTPUT ROUTINE HELPERS
%
% Resolves inner/outer to left/right at output time, based on:
%   \ifdoublesided — whether pages alternate
%   \ifbindingleft — which physical side is the spine
%   \ifodd\pageno  — recto (odd) or verso (even)
%
% On a recto page with binding=left:  inner=left,  outer=right
% On a verso page with binding=left:  inner=right, outer=left
% binding=right simply swaps the above.
%
% \ifinnerisleft is set true whenever inner resolves to left,
% and used by both the headline/footline and the margin note macro.
% =============================================================

\newif\ifinnerisleft

\def\resolvebinding{%
  \ifdoublesided
    \ifodd\pageno
      % recto page
      \ifbindingleft \innerislefttrue \else \innerisleftfalse \fi
    \else
      % verso page
      \ifbindingleft \innerisleftfalse \else \innerislefttrue \fi
    \fi
  \else
    % single-sided: inner=left always
    \innerislefttrue
  \fi
}

% \makeline{LEFT-TOKS}{CENTRE-TOKS}{RIGHT-TOKS}
% Builds a line with left, centred, and right content.
\def\makeline#1#2#3{#1\hfil#2\hfil#3}

% \resolvedheadline and \resolvedfootline call \resolvebinding
% then wire inner/outer to left/right accordingly:
\def\resolvedheadline{%
  \resolvebinding
  \ifinnerisleft
    \makeline{\the\headinner}{\the\headmiddle}{\the\headouter}%
  \else
    \makeline{\the\headouter}{\the\headmiddle}{\the\headinner}%
  \fi
}

\def\resolvedfootline{%
  \resolvebinding
  \ifinnerisleft
    \makeline{\the\footinner}{\the\footmiddle}{\the\footouter}%
  \else
    \makeline{\the\footouter}{\the\footmiddle}{\the\footinner}%
  \fi
}

\headline={\resolvedheadline}
\footline={\resolvedfootline}

% =============================================================
% MARGIN NOTE MACRO
%
% \marginnote{REF}{NOTE}
%   REF  — user-supplied reference symbol, placed inline and in margin
%   NOTE — note text, deposited in the margin via \vadjust + \llap or \rlap
%
% \vadjust injects a zero-height \vbox at the current line.
% Inside that box, \llap (or \rlap for inner-side notes) shifts
% the note content into the margin column.
%
% The correct side is resolved at typesetting time via \resolvebinding,
% so double-sided documents automatically alternate the margin note side.
%
% \llap  = hang left  of the current position (for outer-side notes
%          when inner is on the left, i.e. note goes to the right —
%          we shift right using \rlap in that case; see below)
%
% For a note on the RIGHT margin:
%   \rlap{\kern\hsize\kern\marginparsep \vbox{...}}
% For a note on the LEFT margin:
%   \llap{\vbox{...}\kern\marginparsep}
% =============================================================

\def\marginnote#1#2{%
  $^{#1}$%                                         % inline reference mark
  \vadjust{%
    \vbox to 0pt{%
      \resolvebinding
      % Determine whether the note goes left or right:
      % notesside=outer and inner-is-left  => note is on the RIGHT
      % notesside=outer and inner-is-right => note is on the LEFT
      % notesside=inner and inner-is-left  => note is on the LEFT
      % notesside=inner and inner-is-right => note is on the RIGHT
      \ifnotesouter
        \ifinnerisleft \noteonright \else \noteonleft \fi
      \else
        \ifinnerisleft \noteonleft \else \noteonright \fi
      \fi
      \vss
    }%
  }%
}

% Place note content in the RIGHT margin:
\def\noteonright{%
  \hbox{\kern\hsize\kern\marginparsep
    \vbox{\hsize=\marginwidth \eightpoint $^{\currentref}$\,\currentnote}%
  }%
}

% Place note content in the LEFT margin:
\def\noteonleft{%
  \hbox{\llap{%
    \vbox{\hsize=\marginwidth \eightpoint $^{\currentref}$\,\currentnote}%
    \kern\marginparsep
  }}%
}

% We need to pass REF and NOTE through to \noteonright/\noteonleft.
% Redefine \marginnote to store them first:
\def\marginnote#1#2{%
  \def\currentref{#1}%
  \def\currentnote{#2}%
  $^{#1}$%
  \vadjust{%
    \vbox to 0pt{%
      \resolvebinding
      \ifnotesouter
        \ifinnerisleft \noteonright \else \noteonleft \fi
      \else
        \ifinnerisleft \noteonleft \else \noteonright \fi
      \fi
      \vss
    }%
  }%
}

% Smaller font for margin notes (\eightpoint is standard in plain.tex):
% If not available, define it:
% \def\eightpoint{\font\eightrm=cmr8 \eightrm}

% =============================================================
% DOCUMENT
% =============================================================

% --- Set flags -----------------------------------------------
\bindinglefttrue
\doublesidedtrue
\notesoutertrue

% --- Set head/foot content -----------------------------------
\headinner={\folio}
\headmiddle={}
\headouter={My Document}

% --- Body text -----------------------------------------------
Lorem ipsum dolor sit amet, consectetur adipiscing elit.%
\marginnote{$\star$}{This note is keyed to the star in the body text.}
Sed do eiusmod tempor incididunt ut labore et dolore magna aliqua.

\smallskip
Ut enim ad minim veniam, quis nostrud exercitation ullamco laboris.%
\marginnote{\dag}{A dagger marks this second note.}
Duis aute irure dolor in reprehenderit in voluptate velit esse cillum.

\bye
